% !TEX program = lualatex
% Auto-generated from YAML (v3) — 12: デジタル化の基礎：メディア / 数値表現

\documentclass[aspectratio=43,professionalfonts,handout,10pt]{beamer}
\usepackage{beamer_template}

\newcommand{\LectureNum}{12}
\newcommand{\LectureTitle}{デジタル化の基礎：メディア / 数値表現}
\newcommand{\CourseName}{情報リテラシー}
\newcommand{\TextPages}{12-14}

\begin{document}
% --- Slide 1: title ---

\begin{frame}{\CourseName\quad 第\LectureNum 回「\LectureTitle 」}
  {\small \TermName\quad \TeacherName\quad ／\quad 教科書 \TextPages}

  \vfill

  \begin{block}{今日の流れ}
    \begin{enumerate}
      \item コミュニケーションと技術（Theory 12）
      \item 情報メディアの起源・発展・メディアリテラシー（Theory 12）
      \item アナログとデジタル（Theory 13）
      \item 情報の単位・2進法・16進法（Theory 13/14）
    \end{enumerate}
  \end{block}
  \vfill

  \begin{block}{今日の目標}
    \begin{itemize}
      \item コミュニケーションの仕組みとメディアリテラシーを理解し，アナログ/デジタルの違いと2進法の基本を説明できる。
    \end{itemize}
  \end{block}
  \vfill

  \begin{exampleblock}{キーワード}
    \small \textbf{コミュニケーション} ／ \textbf{符号化} ／ \textbf{復号} ／ \textbf{メディアリテラシー} ／ \textbf{アナログデータ} ／ \textbf{デジタルデータ} ／ \textbf{ビット（bit）} ／ \textbf{バイト（B）} ／ \textbf{2進法} ／ \textbf{16進法}
  \end{exampleblock}
\end{frame}

% --- Slide 2: コミュニケーションの仕組み (Theory 12 p.34) ---

\begin{frame}{コミュニケーションの仕組み}
  \begin{block}{ポイント}
    \begin{itemize}
      \item コミュニケーションでは，伝え手が情報を符号化してメディアを通じて送り出す
      \item 受け手は受け取った記号を経験・知識をもとに復号して意味を理解する
      \item 符号化と復号がうまくかみ合わないと，円滑なコミュニケーションが成立しない
    \end{itemize}
  \end{block}

  \vfill

  \begin{alertblock}{演習}
    コミュニケーションにおける「符号化」と「復号」を，具体例を挙げて説明せよ
  \end{alertblock}

  \begin{slidenote}
    「符号化」＝伝え手が情報を言葉・文字・記号などに変換すること。「復号」＝受け手が受け取った記号を自分の経験や知識で意味に変換すること。例：「ありがとう」という気持ち → 言葉に変換（符号化）→ 相手が「感謝された」と理解（復号）。文化や方言の違いで符号化と復号がかみ合わないと誤解が生じる。【演習解答例】「ありがとう」という感謝を言葉に変換（符号化）し，相手が感謝と受け取る（復号）。
  \end{slidenote}
\end{frame}

% --- Slide 3: 情報メディアの起源と発展・メディアリテラシー (Theory 12 p.34-35) ---

\begin{frame}{情報メディアの起源と発展・メディアリテラシー}
  \begin{block}{ポイント}
    \begin{itemize}
      \item 古代のメディア：のろし・トーキングドラム・ロゼッタストーン・洞窟壁画
      \item 活版印刷の普及→新聞→ラジオ→テレビ→インターネット
      \item インターネットにより，送り手と受け手の区別がなくなり誰もが発信できるようになった
      \item メディアが伝える情報は必ずしも正確・公平ではない
      \item 情報の送り手の意図・立場・目的を意識して読み解き，複数の情報源を比較することが重要
    \end{itemize}
  \end{block}

  \begin{slidenote}
    同じニュースでも，媒体によって見出しや論調が異なる場合がある。SNSでは感情的・断片的な情報が拡散しやすい（フェイクニュース）。情報を受け取るときは「誰が・なぜ・何のために発信しているか」を意識し，複数の情報源で確認することが重要。
  \end{slidenote}
\end{frame}

% --- Slide 4: アナログとデジタル (Theory 13 p.36) ---

\begin{frame}{アナログとデジタル}
  \begin{block}{ポイント}
    \begin{itemize}
      \item \textbf{アナログデータ}：色・温度・時間など，連続的な量で表されるデータ
      \item \textbf{デジタルデータ}：連続的な量を一定の間隔で区切って数値で表したデータ
      \item 自然界の情報はアナログ → コンピュータで扱うにはデジタルに変換が必要
      \item A/D変換（アナログ→デジタル）によってコンピュータで処理できるようになる
    \end{itemize}
  \end{block}

  \vfill

  \begin{alertblock}{演習}
    カメラ・音楽・時計など，アナログからデジタルに変わったものを挙げ，\\
    デジタル化によって何が変わったかグループで話し合おう
  \end{alertblock}

  \begin{slidenote}
    \footnotesize アナログ＝連続量，デジタル＝飛び飛びの数値。例：温度計の目盛り（アナログ）↔ デジタル温度計（デジタル）。CDは音波を44,100回/秒サンプリングしてデジタル化。【演習例】フィルムカメラ→スマホ（消せる），レコード→ストリーミング（いつでも聴ける）。
  \end{slidenote}
\end{frame}

% --- Slide 5: 情報の単位 (Theory 13 p.37) ---

\begin{frame}{情報の単位}
  \centering
  \begin{tabular}{lll}
    \toprule
    単位 & 読み方 & 大きさ \\
    \midrule
    bit & ビット & 情報量の最小単位（0または1） \\
    B & バイト & 1B = 8bit \\
    KB & キロバイト & 1KB = 1,024B \\
    MB & メガバイト & 1MB = 1,024KB \\
    GB & ギガバイト & 1GB = 1,024MB \\
    TB & テラバイト & 1TB = 1,024GB \\
    \bottomrule
  \end{tabular}

  \vfill

  \begin{alertblock}{演習}
    \begin{enumerate}
      \item 5MBは何KBか
      \item 3GBは何MBか
    \end{enumerate}
  \end{alertblock}

  \begin{slidenote}
    「1,024」が基準なのは，コンピュータが2進法を使っているから（$2^{10}=1{,}024$）。スマホ写真1枚≈3〜5MB，4GB USBに写真約800〜1,000枚。通信速度は「Mbps（Mbit/秒）」，ストレージ容量は「MB・GB」で表すことが多い。【演習解答】①$5\times1{,}024=5{,}120$ KB　②$3\times1{,}024=3{,}072$ MB
  \end{slidenote}
\end{frame}

% --- Slide 6: 2進法の基本と変換例 (Theory 14 p.38) ---

\begin{frame}{2進法の基本と変換例}
  \centering
  \begin{tabular}{lll}
    \toprule
    10進法 & 2進法 & 計算（各桁の重み：2³=8, 2²=4, 2¹=2, 2⁰=1） \\
    \midrule
    5 & 0101 & 4×1 + 2×0 + 1×1 = 5 \\
    10 & 1010 & 8×1 + 4×0 + 2×1 + 1×0 = 10 \\
    13 & 1101 & 8×1 + 4×1 + 2×0 + 1×1 = 13 \\
    23 & 10111 & 16+4+2+1 = 23 \\
    \bottomrule
  \end{tabular}

  \vfill

  \begin{alertblock}{演習}
    \begin{enumerate}
      \item $1101_{(2)}$ を10進法で表せ
      \item $23_{(10)}$ を2進法で表せ
    \end{enumerate}
  \end{alertblock}

  \begin{slidenote}
    2進→10進：各桁に重み（$2^0=1,\ 2^1=2,\ 2^2=4,\ 2^3=8,\ \ldots$）を掛けて合計。10進→2進：「2で割り続けて余りを下から読む」。例：$23\div2=11\cdots1,\ 11\div2=5\cdots1,\ 5\div2=2\cdots1,\ 2\div2=1\cdots0,\ 1\div2=0\cdots1$ → 下から読んで $10111$。【演習解答】①$1\times8+1\times4+0\times2+1\times1=13$　②$10111$
  \end{slidenote}
\end{frame}

% --- Slide 7: 16進法 (Theory 14 p.39) ---

\begin{frame}{16進法}
  \begin{block}{ポイント}
    \begin{itemize}
      \item 0〜9とA〜F（A=10, B=11, … F=15）の16文字で表す記数法
      \item 2進数4桁が16進数1桁に対応するため，変換が容易
      \item 例：$11011110_{(2)} \to \mathrm{DE}_{(16)}$（4ビットずつ区切って変換）
      \item 逆に16進→2進：各桁を4ビットの2進数に置き換える
    \end{itemize}
  \end{block}

  \vfill

  \begin{alertblock}{演習}
    $11011110_{(2)}$ を16進法に変換せよ
  \end{alertblock}

  \begin{slidenote}
    2進→16進の変換：右から4ビットずつ区切り，各グループを16進1桁に変換する。例：$\underbrace{1101}_{\text{D}}\underbrace{1110}_{\text{E}}{}_{(2)}=\text{DE}_{(16)}$。Webの色指定 \texttt{\#FF0000} → R=FF(255), G=0, B=0 → 赤。MACアドレス（例：00-1A-2B-3C-4D-5E）も16進表記。【演習解答】1101=D，1110=E → $\mathrm{DE}_{(16)}$
  \end{slidenote}
\end{frame}

% --- チェックリスト ---

\begin{frame}{まとめ・チェックリスト}
  \begin{block}{自己チェック（できたら $\checkmark$ をつけよう）}
    \begin{itemize}
      \item[$\square$] コミュニケーションにおける符号化・復号とメディアリテラシーの重要性を理解した
      \item[$\square$] アナログとデジタルの違い，ビット・バイトの単位を学んだ
      \item[$\square$] 2進法・16進法の仕組みと10進法との変換方法を習得した
    \end{itemize}
  \end{block}

  \vfill

  {\small 質問があれば授業後またはTeamsで受け付けます。}
\end{frame}

\end{document}
